\documentclass[11pt]{article}
\usepackage{geometry}
\usepackage{braket}
\usepackage{amsmath}
\geometry{margin=1in}

\title{Quantencomputer basierend auf photonischen Qubits}
\author{Maximilian Nidens}

\begin{document}
\section{Kriterien von DiVincenzo}
\subsection{Skalierbares physikalisches System mit wohl charakterisierbaren qubits}
Zu beginn benötigt man "einfach" ein quantenmechanisches System mit zwei Zuständen. Einfache Beispiele dafür wären der
Grundzustand und erster angeregter Zustand eines Spin-1/2-Teilchens oder in unserem Fall die vertikale und horizontale
Polarisation eines Photons.\\
Diese werden dann mit $\ket{0}$ und $\ket{1}$ gekennzeichnet. Die wohl wichtigste Unterscheidung von Qubits zu klassischen bits
ist, dass die Zustände eines einzelnen Qubits element eines zwei-dimensionalen komplexen Vektorraums sind. Somit gilt für
einen allgemeinen Zustand
\begin{align}
    a\ket{0} + b\ket{1}
\end{align} 
mit der Normierung $|a|^2+|b|^2 = 1$. Analog folgt für einen allgemeinen zwei-qubit-Zustand
\begin{align}
    a\ket{00} + b\ket{01} + c\ket{10} + d\ket{11}  
    \label{2qubitstate}
\end{align}
und so weiter. Der Zustand in \eqref{2qubitstate} beschreibt dann analog zum ersten Fall einen vier-dimensionalen vektor, wobei
jede Dimension für einen unterscheidbaren Zustand des Systems steht und alle diese Zustände im allgemeinen verschränkt sind.
Das bedeutet, dass sie nicht als Produktzustand der individuellen Qubits schreiben lassen. Verschränkte Zustände lassen somit
eine Rechenpower von $2^n$ zu.\\
Ein wohl charakterisierbarer Qubit bedeutet unter anderem, dass die internen Parameter des Systems, sowie der zugrundeliegende
Hamiltonian bekannt sind. Somit sind überhaupt erst die Energieeigenzustände bekannt und es können Observablen berechnet werden.
Des Weiteren sollte, sofern das System Eigenzustände höherer Energie hat, sichergestellt werden, dass die Übergänge in diese
höheren Energieniveaus unwahrscheinlich sind. Das ist notwendig, damit unsere Qubits weiterhin 2 Level Systeme bleiben. 
Jegliche Übergänge in diese höheren Energieniveaus sind mit Rechenfehlern verbunden.\\
\subsection{Präparation des Zustand zu einem Referenzzustand}
Um überhaupt Berechnungen auf einem Quantencomputer durchzuführen, muss der Ausgangszustand des Systems bekannt sein, da 
das Endergebnis sonst völlig nutzlos wäre, analog zu einer Rechenoperation bei einem klassischen Computer. Üblicherweise ist
dieser Referenzzustand charakerisiert durch $\ket{000...}$. Somit ist der niedrigste Energiezustand des System von großer 
Bedeutung.\\
Dies kann durch mehrere Prozesse geschehen. Zum einen kann das System schlicht gekühlt werden, womit das System auf natürliche
Weise in den Grundzustand übergeht. Eine weitere Mögichkeit wäre, Messungen durchzuführen, welche das System auf den Grundzustand
projizieren oder eine unitäre Operation durchzuführen, welche das System in den Grundzustand versetzt.\\
\subsection{Dekohärenzzeit}
Die Dekohärenzzeit eines Qubits charakerisiert das Verhalten mit der Umgebung. Eine einfache Vorstellung davon wäre die Zeit,
die benötigt wird um einen generischen Zustand $\ket{\psi} = a\ket{0} + b\ket{1}$ in eine Mischung zu transformieren. Dies 
wird durch den quantenmechanischen Dichteoperator $\rho$ beschrieben. Für reine quantenmechanische Zustände gilt $\rho^2 = \rho$
und wir können den Zustand als Zustandsvektor $\ket{\psi}$ schreiben, welcher auch als verschränkter Zustand auftauchen kann. Für reine
Zustände liegt die maximale Information vor, da alle weiteren Zustände mit $e^{i\phi}\ket{\psi}$ denselben Zustand beschreiben.
Für gemischte Zustände lässt sich kein Zustandsvektor, weshalb auch keine Phasenbeziehung zwischen verschiedenen einzelnen Zuständen
(Qubits) besteht. Es ist keine Superposition, noch Verschränkung möglich, weshalb gemischte Zustände für Quantencomputer völlig
nutzlos sind. Wenn nun die Dekohärenzzeit die Zeit angibt, nach welcher ein reiner Zustand in einen gemischten übergeht, ist
es von Vorteil diese möglichst lang zu halten, da sonst zwischen Gatter-Operationen die Information verloren gehen kann.
\subsection{Universelle Menge an Quantengattern}
Um Quantenzustände zu funktionierenden Quantencomputern zu machen braucht man logischerweise auch Gatter-Operationen um 
Quantenalgorithmen technisch überhaupt erst möglich zu machen. Solche werden durch eine Abfolge an unitären Operationen erzeugt
die auf die verschiedenen Qubits wirken.\\
Man könnte jetzt meinen, dass man einen reinen Zustand präpariert und dann verschiedenste unitäre Transformationen anwendet, 
bspw. $U_1 = e^{iH_1t\hbar}$, $U_2 = e^{iH_2t\hbar}$, $U_3 = e^{iH_3t\hbar}$ etc. (wobei $H_i$ die hermiteschen 
Hamiltonoperatoren sind), und das dann in den Zeitintervallen 0 bis $t$, danach $t$ bis $2t$ und so weiter. Das Problem 
dabei ist jedoch, dass während der gesamten Zeit dieser Operationen die Dekohärenzuhr tickt. Des Weiteren gibt es nur 
bestimmte Hamiltonians die sich ohne weiteres an und ausschalten lassen, wobei die meisten von denen 2-Körper-wechselwirkungen
betrachten. Für Qubit-Systeme die größer als zwei sind müssen die höheren Wechselwirkungen meistens durch Abfolgen von zwei-Körper
Wechselwirkungen implementiert werden, was durch sogennante CNOT-Gatter realisiert wird.
\subsection{Qubit-spezifische Messkapazität}
Das letzte wichtige Kriterium wäre, dass nach einer Berechnung auf einem Quantencomputer, das Ergebnis gemessen werden kann.
Dazu muss das System in der Lage sein, nach dieser Berechnung in einen wohldefinierten Zustand überzugehen. Wenn die Dichtematrix
eines Zustands bspw. $\rho = p\ket{0} + (1-p)\ket{1} + \alpha\ket{0}\bra{1} + \alpha^*\ket{1}\bra{0}$ ist, sollte mit einer
Wahrscheinlichkeit von $p$ Zustand "0" und mit der Wahrscheinlichkeit $1-p$ "1" ausgelesen werden. Und das alles unabhängig von 
$\alpha$, anderen Parametern des Systems, ohne andere Qubits zu verändern und auf Einflüsse von anderen Qubits zu reagieren.
Das letzte Kriterium scheint besonders streng zu sein, ist jedoch essentiell für einen funktionierenden Computer im allgemeinen.
\subsection{Was ist ein Qubit?}
$(\ket{0}, \ket{1})$\\
$(\ket{00}, \ket{01}, \ket{10}, \ket{11})$\\
$(\ket{000}, \ket{001}, ..., \ket{111})$\\
\subsection{Kodierung}
Polarisationskodierung
$\ket{H} := \ket{0}$ \& $\ket{V} := \ket{1}$\\
Dual-Rail-Kodierung
$\ket{0}\quad$/$\quad\ket{1}$
\subsection{Hadamard Transformation}
Eine Hadamard Transformation in der Dual-Rail repräsentation überführt ein photonisches Qubit in einem Wellenleiter in eine
Superposition beider Wellenleiter:\\
\begin{align*}
    H\ket{0} = \frac{\ket{0} + \ket{1}}{\sqrt{2}}\\
    H\ket{1} = \frac{\ket{0} - \ket{1}}{\sqrt{2}}
\end{align*}
mit
\begin{align*}
    H = \frac{1}{\sqrt{2}}
    \begin{pmatrix}
    1 & 1 \\
    1 & -1
    \end{pmatrix}
\end{align*}
\subsection{Cross-Kerr}
$\chi^{(3)}\approx$10$^{-22}\,$m$^2\,$V$^{-2}$
\subsection{CNOT-Logik}
$\ket{0}_T$ / $\ket{1}_T$ $\rightarrow$ $\ket{1}_T$ / $\ket{0}_T$
\subsection{CZ-Gate}
2-Qubit-Zustand
\begin{align*}
    \left(\frac{\ket{0} + \ket{1}}{\sqrt{2}}\right)\otimes\left(\frac{\ket{0} + \ket{1}}{\sqrt{2}}\right)\\
    &= \frac{\ket{00} + \ket{01} + \ket{10} + \ket{11}}{2}
\end{align*}
Nach dem CZ-Gate erhält man den Zustand\\
\begin{align*}
    \frac{\ket{00} + \ket{01} + \ket{10} - \ket{11}}{2}
\end{align*}
welcher nicht mehr separabel ist. Damit folgt die VErschränkung beider Qubits.
\subsubsection{NS-Gate}
Zu bearbeitender Zustand
\begin{align*}
    \ket{\psi} = \alpha\ket{0} + \beta\ket{1} + \gamma\ket{2}
\end{align*}
wird zu
\begin{align*}
    \ket{\psi} = \alpha\ket{0} + \beta\ket{1} - \gamma\ket{2}
\end{align*}
\subsection{Photon Bunching}
2 Photonen treffen gleichzeitig auf einen Strahlteiler
\begin{align*}
    \ket{11}\quad\rightarrow\quad\frac{\ket{20} - \ket{02}}{\sqrt{2}}
\end{align*}
Superposition von entweder 2 Photonen oben oder unten. Letzter Strahleiter bewirkt
\begin{align*}
    \ket{\Phi} = \alpha\gamma\ket{0101} + \alpha\delta\ket{0110} + \beta\gamma\ket{1001} -\beta\delta\ket{1010}
\end{align*}
\end{document}